\subsection{2.2 prediction: ringing of the step-edge image}\label{app:pred-ringing}

written 2026-10-23 16:29 IST, before any filtering code for problem 2 was written or run. this file is not edited after the experiment.

\subsubsection*{setup}

s(x, y) is 512 x 512, a square of side 256 and intensity B = 200 centred on a background A = 50. the square covers rows and columns 128 to 383, so its left edge lies between columns 127 and 128, at x\_e = 127.5. it is zero-padded to P = 1024 and filtered with a centred, real, zero-phase H. the profile is a row through the middle of the square (row 255), far from the corners. u is the frequency index measured from the centre of the padded spectrum, so index u means u / P cycles per pixel.

\subsubsection*{the 1-d reduction}

near the middle row and away from the corners the image looks like a function of x only. an image that varies only in x has a fourier transform confined to the line v = 0: if f(x, y) = g(x), then F(u, v) = G(u) * P * delta(v). the filter then acts only through its slice H(u, 0), and the filtered row is the 1-d convolution of g with h\_1(x) = ift\{ H(u, 0) \}.

\subsubsection*{ideal low-pass}

the slice of the ideal filter is H(u, 0) = 1 for \textbar{}u\textbar{} \textless{}= D0, else 0, i.e. a cutoff of f\_c = D0 / P cycles per pixel. in the continuous model

\begin{verbatim}
h_1(x) = 2 f_c sinc(2 f_c x) = sin(2 pi f_c x) / (pi x).
\end{verbatim}

a step of height (B - A) at x\_e has response

\begin{verbatim}
g(x) = A + (B - A) [ 1/2 + (1/pi) Si(2 pi f_c (x - x_e)) ],
Si(z) = integral_0^z sin(t) / t dt.
\end{verbatim}

Si has its global maximum at z = pi, where Si(pi) = 1.851937. so

\textbf{(i) overshoot}

\begin{verbatim}
OS = 100 (max g - B) / (B - A) = 100 (Si(pi)/pi - 1/2) = 8.95 %,
\end{verbatim}

independent of D0 and P (gibbs). with B - A = 150 the peak is at about 200 + 13.4 = 213.4.

\textbf{(ii) distance from the edge to the first overshoot peak}

\begin{verbatim}
2 pi f_c d = pi  ->  d = 1 / (2 f_c) = P / (2 D0) pixels.
\end{verbatim}

\begin{tabular}{lrrrr}
\hline
D0 & 10 & 20 & 40 & 80 \\
\hline
predicted OS (\%) & 8.95 & 8.95 & 8.95 & 8.95 \\
predicted d = P / (2 D0) (px) & 51.2 & 25.6 & 12.8 & 6.4 \\
\hline
\end{tabular}

\subsubsection*{expected sources of discrepancy (also predicted before running)}

\begin{enumerate}
\item \textbf{discrete passband.} on the P grid the ideal slice keeps the 2 D0 + 1 integer frequencies \textbar{}u\textbar{} \textless{}= D0. the kernel is then a dirichlet kernel, h[x] = sin(pi (2 D0 + 1) x / P) / (P sin(pi x / P)), whose first zero is at x = P / (2 D0 + 1). the overshoot peak sits at that zero, so the discrete prediction is d = P / (2 D0 + 1) = 48.8, 25.0, 12.6, 6.4 px. the overshoot of a dirichlet partial sum tends to the same 8.95 \% as the number of terms grows.
\item \textbf{sampling of the profile.} the edge is at x\_e = 127.5 and the profile is only known at integer columns, so a measured d is a half-integer and can be off by up to 0.5 px. the sampled maximum can also miss the true peak, which lowers the measured OS slightly, most for D0 = 80 where the peak is narrow.
\item \textbf{other edges.} the square is finite, and the zero padding adds a jump of -50 at the image border, 128 px from each square edge. for large z, Si(z) \textasciitilde{} pi/2 - cos(z)/z, so the ripple from an edge at distance r has amplitude about P / (2 pi\textasciicircum{}2 D0 r) of that edge's jump. for D0 = 10 this is about 2 \% of the jump for the opposite edge (r = 256) and 1.4 \% for the border jump (r = 128, jump 50). these ripples add or subtract depending on phase, so at D0 = 10 the measured OS can differ from 8.95 \% by a few percentage points, and the difference should shrink roughly like 1 / D0.
\item \textbf{the image is not really 1-d.} the square is only 256 px tall, so its spectrum is spread over a band of v around 0, of width about P / 256 = 4 bins. on the slice at height v the ideal cutoff is sqrt(D0\textasciicircum{}2 - v\textasciicircum{}2) \textless{} D0, which matters most at D0 = 10. the top and bottom edges also send their own ripple down to the middle row. both effects again fade as D0 grows.
\item \textbf{near a corner} the 1-d reduction fails completely: the local image varies in both x and y, its spectrum is not on the v = 0 line, and the 2-d ideal kernel is a jinc whose rings are circular. ripples from the vertical and horizontal edges overlap there, so the profile near a corner should show a different overshoot and a different peak position.
\end{enumerate}

\subsubsection*{gaussian low-pass}

H(u, 0) = exp(-u\textasciicircum{}2 / (2 D0\textasciicircum{}2)). its inverse transform is again a gaussian,

\begin{verbatim}
h_1(x) = exp(-x^2 / (2 sigma_s^2)) / (sqrt(2 pi) sigma_s),
sigma_s = P / (2 pi D0),
\end{verbatim}

which is positive everywhere and integrates to one. the edge response is

\begin{verbatim}
g(x) = A + (B - A) Phi((x - x_e) / sigma_s),
\end{verbatim}

with Phi the standard normal cdf. this is monotone, so

\textbf{predicted gaussian overshoot: OS = 0 \% for every D0.}

the transition just gets wider as D0 drops (10 \% to 90 \% width 2.56 sigma\_s = 2.56 P / (2 pi D0)). on the discrete grid the slice is cut off at \textbar{}u\textbar{} = P/2 = 512, where it is exp(-512\textasciicircum{}2 / (2 D0\textasciicircum{}2)), which is far below float64 precision for D0 \textless{}= 160, so the discrete kernel stays non-negative to round-off and OS should be 0 up to about 1e-10 \%.

\subsubsection*{butterworth (not required, written for 2.2b)}

for n = 1 the slice 1 / (1 + (u/D0)\textasciicircum{}2) is a lorentzian whose inverse transform is a two-sided exponential, positive everywhere, so OS = 0. for n \textgreater{}= 2 the kernel has negative lobes and OS \textgreater{} 0, growing with n toward the ideal-filter limit of 8.95 \% as n -\textgreater{} infinity, since the butterworth slice tends to the ideal box.
